\subsection{Recent Massey theorems and the fixed ordinary problem}

\paragraph{KT-MASSEY-LITERATURE-SCOPE.}
Three further primary-source comparisons were made against the
specific relative obstruction in
KT-ORDINARY-CUBIC-MASSEY-CONDITION. None of the inspected statements
evaluates that obstruction.

Maire--Min\'a\v{c}--Ramakrishna--T\^an,
\emph{On the strong Massey vanishing property for number fields},
Mathematische Annalen 395 (2026), article102,
\url{https://doi.org/10.1007/s00208-026-03518-1},
Theorem1.1, concerns the absolute and maximal tame Galois groups
when $\zeta_p\notin K$. Definition1.1 and the following paragraph
distinguish the existence of one lifting defining system from the
lifting of every prescribed defining system. Remark1.1 explicitly
allows additional ramified primes in the construction; Theorem1.2
retains such primes for higher coefficient lifts. The end of the
introduction fixes trivial coefficients. Our problem has a fixed
finite set $S$, nontrivial character coefficients, and a specified
ordinary defining cochain. Trivializing the coefficient characters
does not supply descent with this fixed ramification set and
ordinary normalization.
Thus the inspected result cannot simply be applied to (CM.6).
The complete statements and these scope paragraphs were read;
the proof of the global embedding theorem was not re-proved.

Ray--Sujatha, \emph{Massey products and the Iwasawa theory of fine
Selmer groups}, author preprint arXiv:2508.17156v1,
\url{https://arxiv.org/pdf/2508.17156v1}, Proposition4.6 on p.14,
relates unobstructedness in an infinite extension to a fine-Selmer
$\mu$ condition under weak Leopoldt. Theorem7.1 on p.24 gives an
equivalence with spanning by products formed from repeated copies
of the extension character $\kappa$ and varying permitted defining
systems. Theorem6.10 and Corollary6.11 on p.23 specify that system.
These are neither evaluation of the fixed alternating-weight
product $(u,v,u)$ nor finite-level vanishing. In addition, the
ordinary application in Section5 assumes weight two and residual
irreducibility over $K(\mu_p)$, explicitly before Proposition5.2.
The residual three-character sum here does not satisfy those
hypotheses. The cited statements, their displayed proofs, and the
Section5 setup were read, without importing an unstated comparison
between the two ordinary deformation problems.

Ahlqvist--Carlson, \emph{Massey products in the \'etale cohomology
of number fields}, Journal f\"ur die reine und angewandte
Mathematik823 (2025),61--112,
\url{https://doi.org/10.1515/crelle-2025-0006},
Corollary4.10 on p.88 and Theorem4.12 on p.90 give explicit
arithmetic formulas on $X=\operatorname{Spec}\mathcal O_K$.
The latter uses an unramified degree-$p$ extension and ideal/norm
data to evaluate a repeated product, modulo its stated
indeterminacy. Its complete statement and short proof were read;
the preceding resolution calculations were not independently
re-proved. In our problem the positive direction may ramify
above $p$, and the local ordinary normalization must be retained.
Remark4.14 on p.93 already sketches the comparison for a punctured
arithmetic curve containing the required roots of unity, referring
to Sharifi's Theorem5.4. That observation must be retained rather
than treating every punctured-curve comparison as absent.
What is still required here is a comparison with the precise
relative ordinary defining system, followed by an actual evaluation
of its norm data. Neither follows merely from identifying a
class-group character.

\paragraph{Arithmetic-geometric location.}
The distinction is between extensions of finite coefficient
sheaves on a fixed punctured arithmetic curve, with local
conditions, and extensions after allowing additional punctures
or after passage to an infinite tower. The finite sheaf setting
is the one already specified at (FW.1)--(FW.2), with SGA4,
Expos\'eIX,3.2--3.3 as its Kummer-theoretic reference. This audit
does not change the remaining arithmetic Triple-NV gate in the
Teichm\"uller-tower crosswalk to the \emph{Esquisse}, Section2,
and the 1983 letter to Faltings, or prove a new integral Tate vanishing.
